% Fichier complété par tools/complete_missing_corrections.py.
% Source : DYN/DYN-03-Inertie/45_Disque/45_Disque.tex
% Statut : rédigé (2 question(s)).
\begin{corrigebox}[Corrigé rédigé]
\CorrigeQuestion{1}
Le centre d'inertie d'un disque homogène est son centre géométrique;
                pour un disque évidé ou composé, on applique le barycentre pondéré en
                considérant les évidements comme des masses négatives.
\CorrigeQuestion{2}
Pour un disque mince de rayon $R$ :
                $I_{zz}=\tfrac12mR^2$ et $I_{xx}=I_{yy}=\tfrac14mR^2$ au centre.
                La matrice en $O$ est obtenue par le théorème de Huygens et sommation des
                différentes parties.
\end{corrigebox}
